% ============================================================
% 6. Cost
%
% GOAL (PAPER_STRUCTURE.md): sparse-memory work is billed per tile touched,
% support clustering rather than support size is the cost variable, and RoLA
% meets attention's FLOP count at N = L.
%
% PROVENANCE RULING (2026-08-02, user). The paper prints only ANALYTIC
% quantities (derived FLOP counts, structural counts, asymptotics) and
% measurements with PUBLISHABLE provenance (rental-class hardware, ratified
% protocol). Every constant below was taken on one consumer card with
% single-launch counters, so NONE of it may be printed: the quanta table's work
% column, the quanta statistics (0.97 / 0.606 / 39 %), the gather A/B
% percentages, the timing ratio, the selection counts and the residuals are all
% DEMOTED to \placeholder / \TODO in the rendered text. The block below is the
% provenance record for the re-measurement, not a source to re-inline from.
%
% NUMBERS. Every constant is READ BACK from the committed cost model,
% `rola/planner/schedule.py` in the standalone package (the file the planner
% selects with), not from the earlier draft and not from memory:
%   GATHER_CANDIDATE_FIXED_WORK   = 177.75    per candidate (owner, tile) pair
%   GATHER_CANDIDATE_QUANTUM_WORK = 525.125   per ceil(BC/16) quantum
%   READOUT_FOLD_KSTEP_WORK       = 1755.0    per k-step
%   INTRA_GRAM_KSTEP_WORK         = 144.0     per k-step
%   DECAY_KSTEP_WORK              = (1106.3 + 959.7) / 2 = 1033.0 per k-step
% The gather is NOT the single ~2278 of the earlier draft: 2278.25 is what the
% BC = 64 cell measures per tile, and the two-point solve over ceil(BC/16)
% against the BC = 16 cell (702.88) splits it into 525.125 per quantum plus a
% 177.75 per-tile intercept. Pricing it as one constant made compaction
% unreachable at BC = 16. Provenance for the solve: rewrite.md Secs. 2.20, 2.33,
% and the constant's own docstring block in schedule.py.
%
% Instrument, all constants: RTX 3080 Ti, `ncu --clock-control base`,
% `smsp__inst_executed.sum`, ONE profiled launch per arm, flock-serialized;
% cells from the frozen registry in `benchmarks/bench_consumer_new.py`.
% Isolating cells: GATHER_* from topo-D2-w64x64-BC64 against topo-D2-w64x16-BC16
% (unions saturate at both, so gathering changes no extent and only the dense
% quantum differs); READOUT_FOLD from atscale-L2048, where extents do move;
% INTRA_GRAM from wk1-D2-w64x64-BC64 against topo-D2-w64x64-BC64, which differ
% only in the intra extent KI; DECAY from the BC = 64 / BC = 16 pair at both
% compact arms, which differ only in the k-step count.
%
% atscale-L2048 config, read back from bench_consumer_new.py: widths (256, 256),
% so N = 65536 at D = 2; B = 1, H = 8, d_v = 64, T = 2048, BC = 64, BT = 16;
% order-1.5 entmax routing PRODUCED by a real RoLA layer (route_W scaled 8.0,
% seed 0) with the plan chosen by the planner. Quanta stats (--mode gatherstats,
% rewrite.md Sec. 2.20): mean K / BC = 0.97, readout k-steps 0.606 of 4, ~39 % of
% candidate pairs gather K = 0. Compaction A/B at that cell: 1812.1 M -> 1035.0 M
% instructions (-42.9 %) at decay = None and 2719.6 M -> 1092.6 M (-59.8 %) with
% decay; paired interleaved timing 8.457 -> 6.383 ms (x0.753, IQR 0.094).
% Dense toll, topo-D2-w64x64-BC64: 9.410 M -> 11.743 M (+24.8 %) and
% 13.888 M -> 15.786 M (+13.7 %).
%
% RESIDUALS (rewrite.md Secs. 2.36-2.37, and the landed test
% `test_the_absolute_instruction_totals_are_a_RECORDING_not_an_assertion`):
% the model reproduces 8 of 8 measured joint (BT, compact) argmins, which is what
% binds the constants; its ABSOLUTE totals are recorded, not asserted, at up to
% +83.9 % with the un-identified per-tile term and -61.0 % (one-sided) without
% it. Decay is overpriced by 2.3x on the one at-scale compact cell (453 measured
% against 1033), stated as a direction rather than fitted with a second constant.
% The earlier draft's "reproduces on a minority of cells, one residual of 2.3x"
% is STALE: that was the pre-calibration state.
%
% The 1.5-2x cost of making both sides dense is NOT quoted as a number: no
% in-tree measurement was found for it, so the rider is stated structurally.
%
% R2 SHAPE (cut plan, 2026-08-03, user-ratified). The section is seven beats,
% law first: (1) the billing law; (2) what the quantum does to an average;
% (3) the two consequences for the architecture; (4) residency in ONE
% paragraph; (5) the FLOP ledger, with the projection arithmetic displayed as
% Eq. (projections) and the adversarial-regime rider and the flat-training-
% memory prediction closing it; (6) the residue and why we keep it; (7) scope
% and the artifact. Compressed here: P6 (residency mechanism), P12 (the
% bench-harness description), P14/P15 (the residue walk-through). Relocated:
% tab:quanta to App. B beside the residual audit that qualifies it. Every cut
% sentence's CLAIM survives, and no cut moved a number into new prose.
% ============================================================
\section{Cost}
\label{sec:cost}

This section prices the head on the hardware that runs it. Accelerators issue work in blocks, so a sparse operator is billed for the blocks it touches and not for the entries it reads. Block-sparse and mixture-of-experts kernels design against that law \citep{chen2022pixelated, guo2025sonicmoe}, and so do we.

What we add is where a routed head lands on it. Think of the state as a grid cut into blocks: a work item pairs an owner, that is, one block of $B_C$ leaf features, with a tile of $B_T$ tokens, and the device stages or skips it whole (Appendix~\ref{apd:costdev}).

The quantization cuts both ways, and no average shows it: a candidate block holds fewer live features than the quantum it is billed for, and other candidate pairs hold nothing and skip every step. \placeholder{At $\Nst = 65{,}536$: mean live features per candidate block of $64$, k-steps paid of four, and the share of candidate pairs that skip; at-scale cell, rental-class hardware, ratified protocol.} A cost model prices the quantum, not the mean.


Two consequences follow, and both invert the intuition that a sparse memory should touch as little as it can. Support size is nearly free once a support has paid for its tiles, so no cost argument remains for a hard per-token budget and Section~\ref{sec:allocation} can let the activation set its own boundary (Appendix~\ref{apd:costdev}).

The second consequence is that clustering, not sparsity, is the quantity worth optimizing: a tile is live if any of its $B_T$ tokens touches it, so the kernel pays for the tokens' combined footprint (Appendix~\ref{apd:costdev}).

Nearby tokens choosing the same features is coherence, that is, tokens in one tile agreeing on where they write, and the factorized address supplies part of it by layout, since a coarse digit selects a contiguous region (Appendix~\ref{apd:costdev}).

What that gives is structural, an address whose contiguity a flat address of the same capacity would have to learn. Hierarchy is a coherence prior in that sense, the cost-side counterpart of the addressing argument of Section~\ref{sec:capacity} (Appendix~\ref{apd:costdev}).

% R2 beat 4 (cut plan, 2026-08-03): the residency CLAIM is that resident bytes
% follow what a sequence wrote, and that translation is amortized to a part in
% a thousand; both stay, with the granule size kept because the "part in a
% thousand" is derived from it. Cut to the artifact: the page table's key, the
% decode walk over the exact support, and the checkpoint serializing only the
% granules a sequence wrote. Those are mechanism a reader can read in the code.
Provisioned capacity and resident memory come apart, which is the same argument on the storage side. The state is paged at a fixed granule and a page is admitted from the realized write support, so resident bytes scale with what a sequence wrote and not with the provisioned count (Appendix~\ref{apd:costdev}).

In effect, a head may provision far more features than the memory it ever occupies. Translation costs a part in a thousand of the state traffic it rides on \placeholder{translation overhead as a share of state traffic at the at-scale cell, rental-class hardware, ratified protocol}.

% R2 (cut plan): tab:quanta RELOCATED to App. B (Limits of the cost evidence),
% where the residual audit that qualifies it already lives. With its work
% column \TODO-demoted the table prints three columns of mechanism names, so
% the body keeps only the sentence naming the terms. \label{tab:quanta} moved
% unchanged; Sec. 8 and fig:surfaces' caption still reference it.
Measurement recovers that law, separating a fixed per-tile term from a per-column one and pricing decay per k-step and not per tile, which Appendix~\ref{apd:costlimits} sets out with the cell that isolates each. Enabling the clock of Section~\ref{sec:decay} is therefore cheap in the regime it targets, since a leaf outside the write support adds no tiles.

\subsection{The FLOP ledger}
\label{sec:ledger}

We price the comparison with softmax attention per pair and then count pairs, since attention's per-token cost grows through the sequence and any one token's price misstates the total. Attention pays $\dqk + \dv$ multiply-accumulates per query-key pair, which is $2\dv$ at the standard head shape $\dqk = \dv$ that the ledger below assumes, and RoLA pays about $\dv$ per (token, leaf) pair (Appendix~\ref{apd:costdev}).

The pair counts run the other way, causal attention visiting the triangle and RoLA the square its dense read sweeps. At $\Nst = L$ the two prefill totals meet.
\begin{align}
  \text{attention} \;\approx\; 2\dv \cdot \tfrac{L^2}{2} \;=\; L^2 \dv,
  \qquad
  \text{RoLA} \;\approx\; \dv \cdot L\Nst \;=\; L^2 \dv .
  \label{eq:parity}
\end{align}
The two factors of two cancel, so the totals are equal.

Two riders bound that equality. The square is the count for the wiring whose read is dense over the leaves. Where a level reads sparsely the tile skipping above takes the RoLA side below it, so Eq.~\eqref{eq:parity} is an upper bound there and an equality here. Both sides also exclude input projections, and by different amounts. Fix $\Nst = L = 65{,}536$ addressed as two levels of width $256$, with $d_{\mathrm{model}} = 1024$. The two excluded terms then read against the same $L^2 \dv$ as
\begin{align}
  \text{router} \;=\; 2 L\, d_{\mathrm{model}} \Slv \;=\; \tfrac{1}{4} L^2 \dv,
  \qquad
  \text{query, key, value} \;=\; 3 L\, d_{\mathrm{model}} \dqk \;=\; \tfrac{3}{64} L^2 \dv .
\end{align}
Attention's projections are therefore three sixteenths of the router's.

Decode returns the same answer. Over $L$ steps the totals are $L^2\dv$ against $L\Nst\dv$, decode being the triangle paid one row at a time, and no phase tilts the ledger (Appendix~\ref{apd:costdev}).

That size is the adversarial regime of Section~\ref{sec:corner}, sparse writes with dense reads, so Eq.~\eqref{eq:parity} says we reach attention's arithmetic while paying to learn what attention gets free (Appendix~\ref{apd:costdev}).

% R2 (cut plan): the harness description (three modes over (N, L), which mode
% is which) is implementation; the CLAIM is flat training memory per token,
% and it keeps its \TODO. The mode row it waits on is unchanged in App. C.
The affine recurrence makes one prediction about memory. Its backward pass reconstructs no history of the state, which Appendix~\ref{apd:proofs} proves from the detached clock, so training memory per token is predicted flat in $L$, not measured \TODO{pending: the mode-axis runs, which wait on the native backward built after the forward freeze}.

% Provenance of the 10^{-4}: an order-of-magnitude estimate of bit-level
% bookkeeping traffic against a multiply-accumulate, from the class of
% arithmetic-vs-bit-traffic ratios, NOT a measured constant and not sourced to
% any run in App. C. The prose says so; if a reviewer wants a number it becomes
% a \placeholder{}. The 4.2 M decisions and the 8 KB bitmap figure that used to
% sit in this paragraph were derivable from the stated B_C, B_T, Nst, L; they
% were CUT by the R2 compression (cut plan, 2026-08-03) as arithmetic a reader
% can redo, not as claims. The claim (a floor of order 1e-4, not zero) stays.
%
% The keep-the-term sentence states the decision and its structural caveat
% (the term scales with Nst*L, the retiring mechanism does not) rather than
% rejecting the removal silently: standing rule "structural waste beats
% today's payoff" says surface, do not decide.
One term in the total stays linear in $\Nst$ whatever a wiring touches, since the kernel decides candidacy once per (owner, tile) pair, $\Nst L / (B_C B_T)$ of those decisions over a prefill. Those are taken as word operations over the branch bitmaps the routing already publishes.

A bit of bookkeeping runs against a multiply-accumulate at a ratio near $10^{-4}$, which we have not measured and state only to an order of magnitude. The slope-zero line of Figure~\ref{fig:surfaces}(b) is therefore predicted to flatten onto a floor of that order and not onto zero.

We keep the term, and Appendix~\ref{apd:quanta} works out why removing it would cost more than it does, together with the structural caveat that the term scales with $\Nst L$ where the mechanism that would retire it stays fixed.

We claim parity at capacity-fair $\Nst = L$, matched on addressing entropy, and not a throughput multiple at long $L$. Any fixed-state model wins that eventually and by construction. Concurrent work already reports large multiples over FlashAttention at $131$K \citep{ghriss2026kata}, a number that measures the regime rather than the architecture.

Memory feed and exposed latency bind this kernel, not arithmetic. A model counting work alone cannot close that gap. We report the law, and we release the implementation. What the evidence behind the model covers, and where it stops, is Appendix~\ref{apd:costlimits}.

\TODO{Artifact URL once minted.}
